\IfFileExists{stacks-project.cls}{%
\documentclass{stacks-project}
}{%
\documentclass{amsart}
}

% For dealing with references we use the comment environment
\usepackage{verbatim}
\newenvironment{reference}{\comment}{\endcomment}
%\newenvironment{reference}{}{}
\newenvironment{slogan}{\comment}{\endcomment}
\newenvironment{history}{\comment}{\endcomment}

% For commutative diagrams we use Xy-pic
\usepackage[all]{xy}

% We use 2cell for 2-commutative diagrams.
\xyoption{2cell}
\UseAllTwocells

% We use multicol for the list of chapters between chapters
\usepackage{multicol}

% This is generally recommended for better output
\usepackage{lmodern}
\usepackage[T1]{fontenc}

% For cross-file-references
\usepackage{xr-hyper}

% Package for hypertext links:
\usepackage{hyperref}

% For any local file, say "hello.tex" you want to link to please
% use \externaldocument[hello-]{hello}
\externaldocument[introduction-]{introduction}
\externaldocument[conventions-]{conventions}
\externaldocument[sets-]{sets}
\externaldocument[categories-]{categories}
\externaldocument[topology-]{topology}
\externaldocument[sheaves-]{sheaves}
\externaldocument[sites-]{sites}
\externaldocument[stacks-]{stacks}
\externaldocument[fields-]{fields}
\externaldocument[algebra-]{algebra}
\externaldocument[brauer-]{brauer}
\externaldocument[homology-]{homology}
\externaldocument[derived-]{derived}
\externaldocument[simplicial-]{simplicial}
\externaldocument[more-algebra-]{more-algebra}
\externaldocument[smoothing-]{smoothing}
\externaldocument[modules-]{modules}
\externaldocument[sites-modules-]{sites-modules}
\externaldocument[injectives-]{injectives}
\externaldocument[cohomology-]{cohomology}
\externaldocument[sites-cohomology-]{sites-cohomology}
\externaldocument[dga-]{dga}
\externaldocument[dpa-]{dpa}
\externaldocument[sdga-]{sdga}
\externaldocument[hypercovering-]{hypercovering}
\externaldocument[schemes-]{schemes}
\externaldocument[constructions-]{constructions}
\externaldocument[properties-]{properties}
\externaldocument[morphisms-]{morphisms}
\externaldocument[coherent-]{coherent}
\externaldocument[divisors-]{divisors}
\externaldocument[limits-]{limits}
\externaldocument[varieties-]{varieties}
\externaldocument[topologies-]{topologies}
\externaldocument[descent-]{descent}
\externaldocument[perfect-]{perfect}
\externaldocument[more-morphisms-]{more-morphisms}
\externaldocument[flat-]{flat}
\externaldocument[groupoids-]{groupoids}
\externaldocument[more-groupoids-]{more-groupoids}
\externaldocument[etale-]{etale}
\externaldocument[chow-]{chow}
\externaldocument[intersection-]{intersection}
\externaldocument[pic-]{pic}
\externaldocument[weil-]{weil}
\externaldocument[adequate-]{adequate}
\externaldocument[dualizing-]{dualizing}
\externaldocument[duality-]{duality}
\externaldocument[discriminant-]{discriminant}
\externaldocument[derham-]{derham}
\externaldocument[local-cohomology-]{local-cohomology}
\externaldocument[algebraization-]{algebraization}
\externaldocument[curves-]{curves}
\externaldocument[resolve-]{resolve}
\externaldocument[models-]{models}
\externaldocument[functors-]{functors}
\externaldocument[equiv-]{equiv}
\externaldocument[pione-]{pione}
\externaldocument[etale-cohomology-]{etale-cohomology}
\externaldocument[proetale-]{proetale}
\externaldocument[relative-cycles-]{relative-cycles}
\externaldocument[more-etale-]{more-etale}
\externaldocument[trace-]{trace}
\externaldocument[crystalline-]{crystalline}
\externaldocument[spaces-]{spaces}
\externaldocument[spaces-properties-]{spaces-properties}
\externaldocument[spaces-morphisms-]{spaces-morphisms}
\externaldocument[decent-spaces-]{decent-spaces}
\externaldocument[spaces-cohomology-]{spaces-cohomology}
\externaldocument[spaces-limits-]{spaces-limits}
\externaldocument[spaces-divisors-]{spaces-divisors}
\externaldocument[spaces-over-fields-]{spaces-over-fields}
\externaldocument[spaces-topologies-]{spaces-topologies}
\externaldocument[spaces-descent-]{spaces-descent}
\externaldocument[spaces-perfect-]{spaces-perfect}
\externaldocument[spaces-more-morphisms-]{spaces-more-morphisms}
\externaldocument[spaces-flat-]{spaces-flat}
\externaldocument[spaces-groupoids-]{spaces-groupoids}
\externaldocument[spaces-more-groupoids-]{spaces-more-groupoids}
\externaldocument[bootstrap-]{bootstrap}
\externaldocument[spaces-pushouts-]{spaces-pushouts}
\externaldocument[spaces-chow-]{spaces-chow}
\externaldocument[groupoids-quotients-]{groupoids-quotients}
\externaldocument[spaces-more-cohomology-]{spaces-more-cohomology}
\externaldocument[spaces-simplicial-]{spaces-simplicial}
\externaldocument[spaces-duality-]{spaces-duality}
\externaldocument[formal-spaces-]{formal-spaces}
\externaldocument[restricted-]{restricted}
\externaldocument[spaces-resolve-]{spaces-resolve}
\externaldocument[formal-defos-]{formal-defos}
\externaldocument[defos-]{defos}
\externaldocument[cotangent-]{cotangent}
\externaldocument[examples-defos-]{examples-defos}
\externaldocument[algebraic-]{algebraic}
\externaldocument[examples-stacks-]{examples-stacks}
\externaldocument[stacks-sheaves-]{stacks-sheaves}
\externaldocument[criteria-]{criteria}
\externaldocument[artin-]{artin}
\externaldocument[quot-]{quot}
\externaldocument[stacks-properties-]{stacks-properties}
\externaldocument[stacks-morphisms-]{stacks-morphisms}
\externaldocument[stacks-limits-]{stacks-limits}
\externaldocument[stacks-cohomology-]{stacks-cohomology}
\externaldocument[stacks-perfect-]{stacks-perfect}
\externaldocument[stacks-introduction-]{stacks-introduction}
\externaldocument[stacks-more-morphisms-]{stacks-more-morphisms}
\externaldocument[stacks-geometry-]{stacks-geometry}
\externaldocument[moduli-]{moduli}
\externaldocument[moduli-curves-]{moduli-curves}
\externaldocument[examples-]{examples}
\externaldocument[exercises-]{exercises}
\externaldocument[guide-]{guide}
\externaldocument[desirables-]{desirables}
\externaldocument[coding-]{coding}
\externaldocument[obsolete-]{obsolete}
\externaldocument[fdl-]{fdl}
\externaldocument[index-]{index}

% Theorem environments.
%
\theoremstyle{plain}
\newtheorem{theorem}[subsection]{Theorem}
\newtheorem{proposition}[subsection]{Proposition}
\newtheorem{lemma}[subsection]{Lemma}

\theoremstyle{definition}
\newtheorem{definition}[subsection]{Definition}
\newtheorem{example}[subsection]{Example}
\newtheorem{exercise}[subsection]{Exercise}
\newtheorem{situation}[subsection]{Situation}

\theoremstyle{remark}
\newtheorem{remark}[subsection]{Remark}
\newtheorem{remarks}[subsection]{Remarks}

\numberwithin{equation}{subsection}

% Macros
%
\def\lim{\mathop{\mathrm{lim}}\nolimits}
\def\colim{\mathop{\mathrm{colim}}\nolimits}
\def\Spec{\mathop{\mathrm{Spec}}}
\def\Hom{\mathop{\mathrm{Hom}}\nolimits}
\def\Ext{\mathop{\mathrm{Ext}}\nolimits}
\def\SheafHom{\mathop{\mathcal{H}\!\mathit{om}}\nolimits}
\def\SheafExt{\mathop{\mathcal{E}\!\mathit{xt}}\nolimits}
\def\Sch{\mathit{Sch}}
\def\Mor{\mathop{\mathrm{Mor}}\nolimits}
\def\Ob{\mathop{\mathrm{Ob}}\nolimits}
\def\Sh{\mathop{\mathit{Sh}}\nolimits}
\def\NL{\mathop{N\!L}\nolimits}
\def\CH{\mathop{\mathrm{CH}}\nolimits}
\def\proetale{{pro\text{-}\acute{e}tale}}
\def\etale{{\acute{e}tale}}
\def\QCoh{\mathit{QCoh}}
\def\Ker{\mathop{\mathrm{Ker}}}
\def\Im{\mathop{\mathrm{Im}}}
\def\Coker{\mathop{\mathrm{Coker}}}
\def\Coim{\mathop{\mathrm{Coim}}}

% Boxtimes
%
\DeclareMathSymbol{\boxtimes}{\mathbin}{AMSa}{"02}

%
% Macros for moduli stacks/spaces
%
\def\QCohstack{\mathcal{QC}\!\mathit{oh}}
\def\Cohstack{\mathcal{C}\!\mathit{oh}}
\def\Spacesstack{\mathcal{S}\!\mathit{paces}}
\def\Quotfunctor{\mathrm{Quot}}
\def\Hilbfunctor{\mathrm{Hilb}}
\def\Curvesstack{\mathcal{C}\!\mathit{urves}}
\def\Polarizedstack{\mathcal{P}\!\mathit{olarized}}
\def\Complexesstack{\mathcal{C}\!\mathit{omplexes}}
% \Pic is the operator that assigns to X its picard group, usage \Pic(X)
% \Picardstack_{X/B} denotes the Picard stack of X over B
% \Picardfunctor_{X/B} denotes the Picard functor of X over B
\def\Pic{\mathop{\mathrm{Pic}}\nolimits}
\def\Picardstack{\mathcal{P}\!\mathit{ic}}
\def\Picardfunctor{\mathrm{Pic}}
\def\Deformationcategory{\mathcal{D}\!\mathit{ef}}
\usepackage{cleveref}
\crefname{originallemma}{original lemma}{original lemmas}
\Crefname{originallemma}{Original lemma}{Original lemmas}
\makeatletter
\newlabel{flat-lemma-compactifyable}{{0A9Z}{source}{}{sourcepointer-1}{}}
\newlabel{flat-lemma-compactifyable@cref}{{[originallemma][0][]0A9Z}{[1][0][]source}{}{}{}}
\makeatother
\crefname{originaltheorem}{original theorem}{original theorems}
\Crefname{originaltheorem}{Original theorem}{Original theorems}
\makeatletter
\newlabel{flat-theorem-nagata}{{0F41}{source}{}{sourcepointer-2}{}}
\newlabel{flat-theorem-nagata@cref}{{[originaltheorem][0][]0F41}{[1][0][]source}{}{}{}}
\makeatother
\crefname{originallemma}{original lemma}{original lemmas}
\Crefname{originallemma}{Original lemma}{Original lemmas}
\makeatletter
\newlabel{lemma-shriek-composition}{{0F7J}{source}{}{sourcepointer-3}{}}
\newlabel{lemma-shriek-composition@cref}{{[originallemma][0][]0F7J}{[1][0][]source}{}{}{}}
\makeatother
\crefname{originallemma}{original lemma}{original lemmas}
\Crefname{originallemma}{Original lemma}{Original lemmas}
\makeatletter
\newlabel{lemma-shriek-proper-and-open}{{0F7B}{source}{}{sourcepointer-4}{}}
\newlabel{lemma-shriek-proper-and-open@cref}{{[originallemma][0][]0F7B}{[1][0][]source}{}{}{}}
\makeatother
\crefname{originallemma}{original lemma}{original lemmas}
\Crefname{originallemma}{Original lemma}{Original lemmas}
\makeatletter
\newlabel{lemma-shriek-proper-and-open-compose}{{0F7C}{source}{}{sourcepointer-5}{}}
\newlabel{lemma-shriek-proper-and-open-compose@cref}{{[originallemma][0][]0F7C}{[1][0][]source}{}{}{}}
\makeatother
\crefname{originallemma}{original lemma}{original lemmas}
\Crefname{originallemma}{Original lemma}{Original lemmas}
\makeatletter
\newlabel{lemma-shriek-proper-and-open-compose-horizontal}{{0F7D}{source}{}{sourcepointer-6}{}}
\newlabel{lemma-shriek-proper-and-open-compose-horizontal@cref}{{[originallemma][0][]0F7D}{[1][0][]source}{}{}{}}
\makeatother
\crefname{originalreference}{original reference}{original references}
\Crefname{originalreference}{Original reference}{Original references}
\makeatletter
\newlabel{remark-going-around}{{0F7E}{source}{}{sourcepointer-7}{}}
\newlabel{remark-going-around@cref}{{[originalreference][0][]0F7E}{[1][0][]source}{}{}{}}
\makeatother
\title{Pseudo-functoriality of compactly supported derived pushforward: commutativity of the triple composition diagram}
\author{The Stacks Project Authors}
\date{stacks-project:more-etale.tex:local-snapshot-20261012; GNU FDL 1.2 or later}
\begin{document}
\maketitle
\setcounter{subsection}{1}
\noindent\textit{Editorial note.} All delivered mathematical text is the author's own, copied from the pinned Stacks Project chapter file more-etale.tex; no statement or proof step was rewritten, reordered or removed. The counted span is the author environment and its own proof environment. References inside the span are handed to the assembler as navigation stubs that carry the canonical tag of the cited tag only; whether each of them is load-bearing is recorded on the stub itself.\par\medskip
\noindent\textit{Editorial note.} Disclosure: two steps are carried by transplant and analogy rather than execution. (1) The identification of the gamma isomorphisms with exchange maps of the displayed rectangles is made by 'arguing exactly as in the proof of Lemma lemma-shriek-composition' without introducing the notation (proof lines 3628-3633); the rectangles are indicated in the displayed diagram. (2) The final interchange gamma\_A after gamma\_C = gamma\_C after gamma\_A is justified by 'the two squares A and C only intersect in one point (similar to the last argument in Remark remark-going-around)' (proof lines 3640-3642); the remark's last argument is itself anchored to the 2-category references in Categories (its statement lines 3039-3043), and the present case is the same interchange situation. No verbatim 'we omit' occurs in the proof.\par\medskip
\noindent\textit{Allowed standard prerequisites (editorial).}\begin{itemize}
\item Vertical and horizontal composition compatibility of exchange maps (Lemmas lemma-shriek-proper-and-open-compose 0F7C and lemma-shriek-proper-and-open-compose-horizontal 0F7D). cited at proof lines 3633-3637. Factor the rectangle exchange maps: gamma\_\{A+B\} = gamma\_B gamma\_A and gamma\_\{B+C\} = gamma\_B gamma\_C.
\item Going-around remark (Remark remark-going-around, canonical tag 0F7E). cited at proof lines 3640-3642. Supplies the interchange argument pattern: exchange maps for squares meeting in one point commute, anchored there to the 2-category references in Categories.
\item Triple compactification (More on Flatness, Theorem flat-theorem-nagata 0F41 with Lemma flat-lemma-compactifyable 0A9Z). cited at proof lines 3608-3613. Provides the compactifications over which the exchange maps of the displayed diagram are formed.
\end{itemize}\medskip
\ifdefined\hypertarget\hypertarget{sourcepointer-1}{}\fi
\noindent\textit{Source reference.} Original lemma 0A9Z: sources/stacks-chow/tags.txt, source lines 12967--12967. Cited at proof lines 3611-3613: compactifiability underlying the triple compactification.\par\medskip
\ifdefined\hypertarget\hypertarget{sourcepointer-2}{}\fi
\noindent\textit{Source reference.} Original theorem 0F41: sources/stacks-chow/tags.txt, source lines 18886--18886. Cited at proof lines 3611-3613: existence of the compactifications of Z over T, Y over compactified Z, and X over compactified Y.\par\medskip
\ifdefined\hypertarget\hypertarget{sourcepointer-3}{}\fi
\noindent\textit{Source reference.} Original lemma 0F7J: sources/stacks-chow/tags.txt, source lines 19009--19009. Cited at statement line 3603 and proof lines 3629-3630: the gamma isomorphisms of the statement, whose identification with exchange maps follows its proof.\par\medskip
\ifdefined\hypertarget\hypertarget{sourcepointer-4}{}\fi
\noindent\textit{Source reference.} Original lemma 0F7B: sources/stacks-chow/tags.txt, source lines 19001--19001. Cited at proof lines 3631-3632: the exchange maps of the rectangles A+B, B+C, A, C that realize the gamma isomorphisms.\par\medskip
\ifdefined\hypertarget\hypertarget{sourcepointer-5}{}\fi
\noindent\textit{Source reference.} Original lemma 0F7C: sources/stacks-chow/tags.txt, source lines 19002--19002. Cited at proof lines 3634-3637: gives gamma\_\{B+C\} = gamma\_B after gamma\_C.\par\medskip
\ifdefined\hypertarget\hypertarget{sourcepointer-6}{}\fi
\noindent\textit{Source reference.} Original lemma 0F7D: sources/stacks-chow/tags.txt, source lines 19003--19003. Cited at proof lines 3634-3637: gives gamma\_\{A+B\} = gamma\_B after gamma\_A.\par\medskip
\ifdefined\hypertarget\hypertarget{sourcepointer-7}{}\fi
\noindent\textit{Source reference.} Original reference 0F7E: sources/stacks-chow/tags.txt, source lines 19004--19004. Cited at proof lines 3640-3642: the interchange argument pattern for exchange maps of squares meeting in one point, anchored there to the Categories references.\par\medskip
\begin{lemma}
\label{lemma-pseudo-functor}
Let $f : X \to Y$, $g : Y \to Z$, $h : Z \to T$ be separated morphisms of
finite type of quasi-compact and quasi-separated schemes. Then
the diagram
$$
\xymatrix{
Rh_! \circ Rg_! \circ Rf_! \ar[r]_{\gamma_C} \ar[d]^{\gamma_A} &
R(h \circ g)_! \circ Rf_! \ar[d]_{\gamma_{A + B}} \\
Rh_! \circ R(g \circ f)_! \ar[r]^{\gamma_{B + C}} &
R(h \circ g \circ f)_!
}
$$
of isomorphisms of Lemma \ref{lemma-shriek-composition} commutes
(for the meaning of the $\gamma$'s see proof).
\end{lemma}

\begin{proof}
To do this we choose a compactification $\overline{Z}$
of $Z$ over $T$, then a compactification $\overline{Y}$ of $Y$ over
$\overline{Z}$, and then a compactification $\overline{X}$ of
$X$ over $\overline{Y}$. This uses
More on Flatness, Theorem \ref{flat-theorem-nagata} and
Lemma \ref{flat-lemma-compactifyable}.
Let $W \subset \overline{Y}$ be the inverse image of $Z$ under
$\overline{Y} \to \overline{Z}$ and let $U \subset V \subset \overline{X}$
be the inverse images of $Y \subset W$ under $\overline{X} \to \overline{Y}$.
This produces the following diagram
$$
\xymatrix{
X \ar[d]_f \ar[r] & U \ar[r] \ar[d] \ar@{}[dr]|A &
V \ar[d] \ar[r] \ar@{}[rd]|B & \overline{X} \ar[d] \\
Y \ar[d]_g \ar[r] & Y \ar[r] \ar[d] & W \ar[r] \ar[d] \ar@{}[rd]|C &
\overline{Y} \ar[d] \\
Z \ar[d]_h \ar[r] & Z \ar[d] \ar[r] & Z \ar[d] \ar[r] & \overline{Z} \ar[d] \\
T \ar[r] & T \ar[r] & T \ar[r] & T
}
$$
Without introducing tons of notation but arguing exactly
as in the proof of Lemma \ref{lemma-shriek-composition}
we see that the maps in the first displayed diagram use the
maps of Lemma \ref{lemma-shriek-proper-and-open} for the rectangles
$A + B$, $B + C$, $A$, and $C$ as indicated in the diagram in
the statement of the lemma. Since by
Lemmas \ref{lemma-shriek-proper-and-open-compose} and
\ref{lemma-shriek-proper-and-open-compose-horizontal}
we have $\gamma_{A + B} = \gamma_B \circ \gamma_A$ and
$\gamma_{B + C} = \gamma_B \circ \gamma_C$ we conclude
that the desired equality holds provided
$\gamma_A \circ \gamma_C = \gamma_C \circ \gamma_A$.
This is true because the two squares $A$ and $C$ only
intersect in one point (similar to the last argument in
Remark \ref{remark-going-around}).
\end{proof}

\end{document}
